\section{Experimental design}\label{sec:design}
The primary comparison is ESH against ECP inside one implementation, with formulation, initialization, incumbent recovery and tree management held fixed. Comparisons with other solvers, exact cone reformulations and pilot ablations provide context; they do not isolate the separator choice.

\subsection{Models, split and chronology}
The controlled suite contains 51 GDPs: 45 two-product technology-allocation models using five cost laws, and six coupled-region models using sums of exponentials. The cost laws are exponential, negative logarithmic, reciprocal, quadratic and trigonometric. All have nonlinear disjunct constraints. Small, medium and large denote $n=3,8,16$ units, respectively, with three alternatives per unit. The allocation models have $3n$ continuous variables and $3n$ binary indicators; region models have $2n+1$ continuous variables and $3n$ indicators. These are synthetic controls, not calibrated applications. \Cref{app:models} gives the complete model equations, parameter laws, witnesses and convexity proofs.

Seed 104729 supplies the 18 pilot controls. Allocation seeds 130363 and 155921 and region seed 130363 supply the 33 held-out controls. There is no rejection sampling or selection by solver outcome. Allocation coefficients coincide across cost laws at a given seed, and unit parameters share prefixes across sizes. Consequently, family, seed and size observations are related. We report descriptive paired comparisons and all family/size strata, without independence-based tests or confidence intervals.

The original 42 controls have supplied standard-cone formulations. The nine trigonometric controls were added before the main study because the original design did not test a function outside the supplied cone translators; their parameters and settings were fixed without observing their performance. This is a statement about the supplied translators, not impossibility of a conic representation. A preliminary continuous-reference pass evaluated all 42 supported roots before the source freeze. Aggregate status, residual and timing information was inspected, without tuning the algorithm or instances to held-out objectives. Thus ``held out'' means predeclared and untuned, not completely unexamined. Reported references come from a separate, fingerprinted post-freeze pass; the preliminary pass remains archived.

The external set is the complete registry of 27 historical models drawn from GDPlib and Pyomo's GDP examples \citep{gdplib,pyomoexamples}. It includes batch processing, positioning, circle selection, constrained layout, farm layout and small literature examples. Related data and distance-metric variants remain separate runs but are not independent applications. Eighteen models have nonlinear disjunct rows, all quadratic; the remaining nonquadratic expressions occur globally. Thus this external set does not test general nonquadratic disjunct geometry. Its expression-based scope classification is given in \cref{app:legacy-scope}; in particular, all six norm-objective layout cases form a nonsmooth stress stratum, including those that solve successfully.

\subsection{Methods and effective settings}\label{sec:settings}
The primary Cartesian schedule has 13 methods on 51 instances (663 runs). Eight configurations cross ESH/ECP, hull/big-M and single/multiple trees, using the implementation of \cref{sec:algorithm}. Both separators share interior initialization and reduced NLP recovery. The default fractional LP cutoff is $10^{-6}$; optional node user cuts use $0.05$. Neither is the residual-calibrated policy proved in \cref{sec:theory}.

The other methods are GDPopt LOA \citep{chen2018}, SHOT big-M, SHOT declared-convex $\eps$-hull \citep{lundell2022}, Gurobi big-M \citep{gurobi}, and SCIP big-M \citep{scip10}. Pyomo's standard big-M and hull transformations are used for these reformulation baselines. The SHOT hull baseline uses Pyomo's default $\eps=10^{-4}$ perspective approximation, with disaggregated initial values $\nu=\lambda x$ to avoid artificial inactive-branch domain errors. It is not the exact closed-cone hull. The declared-convex SHOT option is restricted to the analytically verified generated models. All SHOT runs use Gurobi for the MIP subsolver, integer tolerance $10^{-8}$, linear feasibility tolerance $10^{-6}$, nonlinear feasibility tolerance $10^{-8}$, and disable trust in cached linear constraint values. These settings were chosen during pilots. The legacy schedule replaces SHOT $\eps$-hull by the supported exact conic adapter, again giving 13 methods and 351 runs. Default SHOT convexity classification remains in effect for the legacy big-M runs.

GDPopt and the prototype use Ipopt \citep{ipopt} for NLPs and Gurobi for MIPs. Overall gap targets are relative $10^{-4}$ and absolute $10^{-6}$ where exposed by the benchmark interface; GDPopt additionally uses bound tolerance $10^{-6}$. At these study tolerances, the prototype's inner Gurobi master uses the tighter settings \texttt{MIPGap}$=10^{-6}$ and \texttt{MIPGapAbs}$=10^{-8}$, shared within each ESH/ECP pair. All native options, mapped statuses and raw bounds are retained. The exact quadratic mixed-integer hull baseline uses Gurobi with \texttt{NonConvex=0} and both \texttt{FeasibilityTol} and \texttt{IntFeasTol} set to $10^{-8}$; \cref{app:cones} gives its formulation and supported scope. It is run on all nine quadratic controls and supported external models. The separate CVXPY/Clarabel reference \citep{diamond2016,goulart2024} solves continuous hull roots and exhaustive small fixed assignments, not general mixed-integer exponential-cone problems. Each reference solve uses one thread, a 300-second limit, at most 300 iterations, and absolute-gap, relative-gap and feasibility targets of $10^{-9}$. These reference limits differ from the benchmark limits below.

The environment is Python 3.13.11, Pyomo 6.10.1, gurobipy/Gurobi 13.0.3 and Ipopt 3.14.20. GAMS 54.3.1 supplies Gurobi 13.0.2, SCIP 10.0.3 and SHOT 1.1 (commit \texttt{a81275b4}); GAMS/SCIP's internal Ipopt is 3.14.19. The separately pinned reference environment uses CVXPY 1.7.3 and Clarabel 0.11.1. Historical solver logs identify an Intel Xeon w5-2565X under WSL2 Linux, with 36 logical CPUs and approximately 30 GiB memory. Runs use one solver thread and one BLAS/OpenMP thread, with at most six concurrent experiment workers; other machine work was present. Each benchmark run has a 120-second solver limit and a 150-second whole-process cap. End-to-end wall time includes interpreter startup, construction, preprocessing, interior solves and optimization.

Primary job order uses seed 20260919. Two complete repetitions of the four single-tree configurations on all 33 held-out controls use order seeds 20260920 and 20260921 (264 additional runs). Solver search seeds are unchanged. Repetitions therefore measure scheduling/runtime variation on the same instances, not search-seed robustness. Separate fixed schedules contain nine quadratic cone runs, 351 external runs and 144 pilot ablations. The ablations cross all 18 pilots and four single-tree configurations with either disabled integer-point NLP recovery or enabled fractional user cuts. Two outcome-triggered follow-ups add nine trigonometric tolerance runs and 24 legacy initialization runs; their declarations and results remain separate. The complete benchmark has $663+264+9+351+144+9+24=1464$ records. References add 42 continuous roots and $14\cdot27=378$ fixed assignments, or 420 solver calls. No repeat or follow-up replaces an original outcome.

\subsection{Acceptance, audit and summaries}\label{sec:acceptance}
A saved point is evaluated on a fresh original GDP, independently of the solver's transformed model. All original numerical variables must be supplied and finite, including otherwise unused variables. Variable bounds, fixed values, active global and selected disjunct rows, integrality, exclusive selection, Boolean consistency and logical constraints are checked. A scalar row $b\le\beta$ passes when
\[
 (b-\beta)_+\le10^{-6}+10^{-7}\max\{1,|b|,|\beta|\};
\]
lower bounds use $(\beta-b)_+$ and equalities require both directions. The same scaled rule applies to variable bounds, with the variable value in place of $b$. Integer distance to the nearest integer must not exceed $10^{-6}$. The original objective is recomputed and checked against a supplied reported objective using the checker tolerance.

An accepted numerical solve requires completed worker output, a valid primal witness, a finite usable solver-reported global bound, and gap
\begin{equation}\label{eq:accept-gap}
 |f(x)-B|\le 10^{-6}+10^{-4}\max\{1,|f(x)|\},
\end{equation}
with the bound in the correct objective sense. Bounds and infeasibility reports are also compared against every validated witness for the same instance across methods, repetitions and feasible reference assignments. Define the objective comparison tolerance at a witness objective $w$ by
\begin{equation}\label{eq:comparison-tolerance}
 T(w)=10^{-6}+10^{-4}\max\{1,|w|\}.
\end{equation}
For minimization, a bound contradicts that witness if $B>w+T(w)$; for maximization, if $B<w-T(w)$. An infeasibility report contradicts the existence of any validated feasible witness. Either contradiction disqualifies an outcome. For GAMS, the unchanged interface uses finite native \texttt{OBJEST} and accepted global-status codes; missing bounds remain missing. An ``optimal'' status alone never establishes acceptance. These checks are numerical evidence, not exact arithmetic certificates.

The final audit checks all saved schedules, unique method--instance pairs, frozen executable/model hashes and environment fingerprints, then re-evaluates supplied witnesses. A separately implemented arithmetic audit, which does not import the study analyzer, reconciles all 1464 benchmark records and 6122 exported analysis fields. It reuses the original-model checker for fresh feasibility evaluation; this is independent aggregation, not an independent second feasibility implementation. All 174 feasible fixed-assignment reference witnesses pass fresh checks, and no cross-witness bound/status contradiction is found. Invalid witnesses, exceptions, missing witnesses, unsupported inputs and timeouts remain visible.

For a cohort $C$ with wall times $t_i$, we use the one-second shifted geometric mean
\[
 S(C)=\exp\!\left(|C|^{-1}\sum_{i\in C}\log(1+t_i)\right)-1.
\]
Paired ratios divide the two shifted means on exactly the same common-solved instances. They are conditional timing comparisons, not failure-penalized speedups. The full-schedule PAR10 arithmetic mean assigns $\min(t_i,150)$ seconds to accepted solves and 1500 seconds to every unaccepted result. A single failure therefore contributes $1500/33\approx45.45$ seconds to a held-out mean. Every table states its denominator. Work means and medians use the same paired cohorts; overlapping component times are never added.
